% Volume IV: Law as an Epistemic Machine
% Part VIII. Legal Friction as Freedom
% Chapter 44: Due Process as Constrained Updating
% Spine: proof-spines/ch044-spine.md

\chapter{Due Process as Constrained Updating}
\label{ch:ch044}

Institutional belief has state \(\Theta_t\). Unrestricted reaction to new evidence \(E_t\) would permit
\[
\Theta_{t+1} = \arg\min_\Theta L(E_t;\Theta),
\]
\ToyModel\ even when \(E_t\) is narrow or unrepresentative.
Due process adds a regularization term:
\[
\Theta_{t+1} = \arg\min_\Theta \left[ L(E_t;\Theta) + \lambda D(\Theta,\Theta_t) + \gamma C(\Theta) \right],
\]
\ToyModel\ where \(D\) penalizes unjustified institutional drift and \(C\) encodes procedural constraints.

Due process does not merely preserve an established prior, since precedent itself can encode systematic error. Its deeper function is to constrain the update rule: new evidence may revise institutional judgment, but it cannot immediately rewrite the entire structure by which evidence is interpreted.

A useful technical analogue appears in parameter-efficient machine learning. In AFF-Net, unrestricted fine-tuning of a large pretrained model on scarce medical data produces feature drift: the model reorganizes broadly acquired representations around the accidental regularities of a small sample and can perform worse than a system permitted to update only through constrained adapters (see Chapter~\ref{ch:ch086} for the original technical context, and Chapter~\ref{ch:ch074} for the full development of the update-budget formalism this chapter introduces informally). Due process imposes an analogous restriction on institutional belief revision. An early recording, witness statement, or accusation may contain genuine information, but it is not statistically entitled to overwrite the entire interpretive structure through which responsibility is assigned. Standards of evidence, cross-examination, precedent, disclosure, and appeal force the new signal to survive contact with a larger body of accumulated constraints. Their value is not that the inherited structure is infallible. It is that limited evidence should initially produce limited revision. The legal equivalent of feature drift occurs when a vivid but narrow sample reorganizes judgment before its provenance, representativeness, and causal context have been tested.

% TODO: the discussion must also address the status-quo-bias
% objection logged in proof-spines/ch044-spine.md — that
% "constrain the update rule" risks sounding like "protect bad
% priors more slowly." The response (already drafted in the
% spine) is that update magnitude is proportional to tested
% reach of new evidence, not direction: a well-tested,
% high-reach body of new evidence should and does produce large
% revision, including overturning precedent. What is restricted
% is a single, untested, high-salience sample producing a large
% revision before its representativeness is known.

The permitted magnitude of an institutional update should be proportionate not to the vividness of the evidence, but to its tested evidentiary reach. The same constraint becomes immunity for bad priors when institutions treat inherited categories as presumptively self-vindicating and keep demanding ever more proof against patterns those categories were built not to see.
